\documentclass{article}
\usepackage{pathfinder-readable}

\title{What Predictions Reveal About Distributed Tail Risk}

\begin{document}
\maketitle

\begin{abstract}
This note reads a survey of learning with predictions alongside a paper on distributed optimisation of
conditional value at risk (CVaR). The thread asked whether predicted loss distributions could reduce the
number of sampled losses needed by the distributed method while remaining safe when predictions are wrong. The
peers obtained an error bound for validated predictions and a lower bound showing that even exact quantile
advice can leave costly tail information hidden. They also recovered a one-query CVaR method, but found that
its centralised estimator and rate were already published. The thread paused because it had neither a
prediction-sensitive query saving nor a proved new distributed result.
\end{abstract}

\section{Introduction}

Q surveys algorithms that use fallible predictions while retaining formal guarantees \cite{Q}. It separates
the form of a prediction from its error, the quality of a result from the cost of obtaining advice, and proved
guarantees from evidence about systems. Its discussion of predicted distributions cautions that a guarantee
for one objective need not carry over to a tail-risk objective. It asks that predictor calls be counted when
comparing methods.

P studies agents that jointly choose a decision while communicating over a changing network \cite{P}. Each
agent sees noisy samples of its own loss, rather than its risk function or gradient. The objective is the
average of local CVaR values, which measure losses in an upper tail. P estimates each CVaR from a fresh batch,
uses a random perturbation to estimate a gradient, and proves an expected bound for an average of iterates.
Its bound records both disagreement between agents and error from finite batches. With fixed batches and
smoothing, its last-iterate result retains a gap from the true CVaR optimum.

The connexion was a concrete cost question: could a forecast of local losses replace some of P's samples, with
a bound that improves when the forecast is accurate and stays valid when it is wrong? The peers treated a
sampled loss and a predictor call as costs, as Q requires. P supplied the oracle and error terms against which
an apparent saving could be tested.

\section{What was done}

The peers first considered a predicted cumulative distribution function (CDF) at each sampled point. They
compared it with an empirical CDF from fresh losses and accepted the prediction only when the two were close.
A concentration bound then controlled the selected CDF, even for a bad forecast. P's inequality between CDF
error and CVaR error turned this into an error-sensitive CVaR estimate. The validation batch still had to be
collected, so this calculation did not save loss queries.

Next, the peers tried subtracting a predicted scalar CVaR value from P's one-point gradient estimate. A scalar
fixed before the random direction has zero mean after multiplication by that direction. This kept the
estimator's conditional mean unchanged and could reduce its second moment when the scalar was accurate. Ada
checked the required timing and the moment bound in \texttt{ada/research-note.md}. A second independent batch
could supply a similar baseline without a prediction, at roughly twice the loss-query cost. Earlier
residual-feedback work also weakened any novelty claim \cite{Residual2010,Wang2022Game}. A forecast benefit
would therefore need a full comparison of charged predictor and loss calls, including P's network error. That
comparison was not completed.

The work then shifted to the threshold used to define CVaR. Rather than estimating CVaR anew from a batch at
every point, the peers kept each agent's threshold as a variable to be optimised. An early update combined a
gradient from ball smoothing with a threshold derivative from sphere smoothing. These are derivatives of
different objectives, as Ada objected. Emmy repaired this by perturbing the decision and threshold together.
Ada checked the repair; their derivations are in \texttt{emmy/lifted-threshold.md} and
\texttt{ada/lifted-threshold.md}. They then found that Cardoso and Xu had already used the same joint
estimator and its centralised rate \cite{CardosoXu2019}.

Finally, the peers asked what an exact threshold forecast actually reveals. They built two loss models with
the same exact quantile at every decision and the same full loss CDF at one anchor decision, but opposite best
decisions. Their information bound quantified the samples still needed to tell the models apart. A separate
suggestion used a shrinking perturbation and P's graph mixing to seek exact last-iterate convergence. The
potential argument in \texttt{emmy/vanishing-radius.md} remained a sketch, and a comparison with existing
distributed one-point work was incomplete \cite{LiAssaad2021}.

\section{Results and their support}

For the validation calculation, let \(F\) be the true loss CDF at a query point, \(\widetilde F\) a predicted
CDF fixed before sampling, and \(F_s\) the empirical CDF of \(s\) fresh losses. There are \(m\) agents and a
horizon of \(T\) rounds. Write \(\eta=\|\widetilde F-F\|_\infty\) for prediction error and set
\(b=\sqrt{\log(2mT/\zeta)/(2s)}\), where \(\zeta\) is the allowed failure probability. The peers' conditional
concentration and union-bound argument gave \(\|F_s-F\|_\infty\le b\) at all the tested points with
probability at least \(1-\zeta\). They selected \(\widetilde F\) when \(\|\widetilde F-F_s\|_\infty\le2b\),
and \(F_s\) otherwise. On that event, the selected CDF has error at most \(3\min\{\eta,b\}\): a forecast with
\(\eta\le b\) is accepted, while any accepted forecast with larger error is within \(3b\) of \(F\). For agent
\(i\), P's CDF inequality then gives a CVaR-value error at most \((6U_i/\alpha_i)\min\{\eta,b\}\), where
\(U_i\) bounds the absolute loss and \(\alpha_i\) is its upper-tail probability. This is a value-error
calculation from the peer record, not a proved query saving or a new network theorem.

The lifted construction explains the alternative sampling route. Let \(x\) be the shared decision,
\(J^i(x,\xi)\) a sampled loss for agent \(i\), and \(\tau_i\) its threshold. Here \(\xi\) denotes a fresh
random sample, and \((a)_+=\max\{a,0\}\). P's CVaR representation uses
\[
H_i(x,\tau_i,\xi)
=\tau_i+\alpha_i^{-1}\bigl(J^i(x,\xi)-\tau_i\bigr)_+,
\qquad
C^i(x)=\min_{\tau_i}\mathbb E H_i(x,\tau_i,\xi).
\]
Under P's bounded, convex loss assumptions, \(H_i\) is jointly convex in \((x,\tau_i)\), and a best threshold
lies in \([-U_i,U_i]\). If \(x\) has dimension \(d\), a random direction \(u\) on the sphere in dimension
\(d+1\) perturbs the pair \((x,\tau_i)\) by radius \(\delta\). One loss sample evaluates \(H_i\) at that pair.
Multiplying the value by \((d+1)u/\delta\) gives a vector whose conditional mean is the gradient of one
jointly ball-smoothed objective. Bounded losses and projection of the stored threshold give a pathwise
estimator bound. The peers' centralised calculation therefore obtained an expected \(O(T^{-1/4})\) gap from
one loss query per agent per round, without P's fixed-batch CVaR-estimation term. Cardoso and Xu already give
this estimator and centralised rate \cite{CardosoXu2019}; the peer calculation checks its fit to P's
formulation. The proposed extension to P's changing graph was not proved.

The lower bound is the clearest limit on advice. For one agent let \(x\in[-1,1]\), let \(E\) equal one with
probability \(\alpha\), and let an independent sign \(S\) equal \(+1\) with probability \(1/2+\gamma\) in one
world and \(1/2-\gamma\) in the other, where \(0<\gamma\le1/4\). A query returns
\[
J(x,E,S)=E(1+Sx)/2.
\]
Both worlds have the exact lower \((1-\alpha)\)-quantile zero at every \(x\). At \(x=0\), both also have the
same full loss CDF: probability \(1-\alpha\) at zero and \(\alpha\) at \(1/2\). Yet their CVaR functions are
\(C_\pm(x)=1/2\pm\gamma x\), so their best decisions are opposite. For \(n\) adaptive loss queries, the peers
bounded the divergence between the two observation histories by \(16n\alpha\gamma^2\): only a tail draw can
reveal the biased sign. Their total-variation argument then gave a worst-world expected CVaR gap of at least
\(\gamma/2\) when \(n\le1/(32\alpha\gamma^2)\). Taking \(\gamma=2\varepsilon\), for \(0<\varepsilon\le1/8\),
yields a necessary order of \(1/(\alpha\varepsilon^2)\) queries for expected gap below \(\varepsilon\). Ada's
proof and Emmy's check are recorded in \texttt{ada/lifted-threshold.md} and the peer ledger. This bound
concerns exact threshold advice, even with an exact CDF at one decision. It does not cover predictions of the
loss distribution across all decisions.

\section{Objections and limits}

The forecast checks separated mathematical validity from a useful saving. Validation spent the batch it was
meant to replace. Scalar centring preserved the conditional mean only when the hint was fixed before the
random direction; a prediction-sensitive network term would require a new disagreement argument. The initial
query-exponent comparison for centring also used an untuned step-size schedule, which Ada corrected. None of
these calculations establishes that predictions beat a fresh baseline batch after all costs are counted. The
ledger also recorded possible one-point and directional-hint comparators without completing the comparison
\cite{Residual2010,DirectionalHint2026}. A separate risk-averse momentum method was logged as relevant prior
work, not as a settled comparison \cite{Wang2022Momentum}.

For threshold lifting, the first derivative pairing failed because it mixed two smoothings. Joint perturbation
settled that proof objection, but the direct Cardoso--Xu precedent settled the centralised novelty claim. The
proposed changing-network bound and the shrinking-radius last-iterate result still need complete proofs and
prior-art checks. Li and Assaad already cover general one-point distributed optimisation on changing networks
\cite{LiAssaad2021}. The two-world lower bound leaves a wide gap against the known \(O(\varepsilon^{-4})\)
upper route and says nothing about richer predictions that reveal how tail losses respond to \(x\).

\section{Why the thread stopped}

The verifier paused the thread because the sound two-world bound limits only a narrow advice model, CDF
validation still uses fresh samples, and the one-query centralised route was already known. The record
contains no prediction-sensitive loss-query saving and no proved new theorem for P's network. The smallest
useful restart is to specify one prediction of tail response across decisions, charge its acquisition and use,
and prove an error-dependent query bound against the one-query lifted baseline and the two-world instance. The
changing-network proof becomes a separate next step if that single-agent comparison succeeds.

\bibliographystyle{plainurl}
\bibliography{references}
\end{document}
